% \newpage
\section{Self-Questioning}
\label{sec:self-questioning}

\begin{figure}
  \centering
  \includegraphics[width=1.0\textwidth]{figs/questioning-overview.pdf}
  \caption{The pipeline of \texttt{self-questioning} module, including \textit{exploration}, \textit{task synthesis}, and \textit{task curation}. Initially, the module performs an exploratory phase across environments, which is directed by both environment profiles and user preferences (detailed in Section \ref{sec:exploration-with-profiles}). The generated trajectories are analyzed to formulate potential queries alongside their associated reference solutions. Next, we involve an agent that replays this path to verify the feasibility.}
  \label{fig:overview-cstb}
\end{figure}

In agent self-evolution, the adaptation process can be characterized as policy optimization over trajectories under proxy-generated tasks and rewards. Consequently, the quality of the self-synthesized tasks directly influences the effectiveness of policy improvement and its eventual transfer to the true target tasks.
However, in practical settings, acquiring training tasks of high quality remains challenging \citep{kang2025demystifying, villalobos2024position}. The difficulties mainly arise from two factors: i) Absence of training data in new environments: Most environments lack predefined data suitable for training, and manually constructing multi-step, complex tasks for agent learning is costly \citep{tang2025synthesizing}. ii) Low quality of training data: Available data frequently suffer from restricted diversity and skewed distributions, which constrain the agent's capacity to develop general and robust capabilities \citep{chen2024diversity}.


To address these challenges and enable agents to continually self-optimize across diverse environments, we propose a \texttt{self-questioning} mechanism
, which serves as the instantiation of the proxy task generation function $F_{\mathrm{task}}$ and synthetic rewards. As shown in Figure~\ref{fig:overview-cstb}, it implements a general exploration-generation pipeline to transform the sandbox $\mathcal E$ into exploration trajectories, and ultimately into user-preferred tasks, leveraging them to bootstrap self-evolution. The \texttt{self-questioning} module tries to answer four main problems in data generation:
\begin{enumerate}
    \item How to understand an unknown environment and find valuable states in it?  (Section \ref{sec:exploration-with-profiles})
    \item How to generate the tasks aligned with user preferences while maintaining sufficient diversity to cover heterogeneous demands? (Section \ref{sec:adaptive-task-synthesis})
    \item How to avoid hallucinations and guarantee the task executability? (Section \ref{sec:task-curation})
    \item How to provide proxy reward signals that evaluate agent performance at the trajectory level? (Section \ref{sec:self-questioning-llm-judge})
\end{enumerate}
By tackling these problems, the proposed module facilitates the co-evolution of tasks and agents, thereby progressively unlocking the agent's potential.

\subsection{Curiosity-Guided Exploration with Environment Profiles} \label{sec:exploration-with-profiles}

To effectively bootstrap synthetic tasks and enable agents to evolve corresponding abilities in a new environment, it is crucial to first develop a structured understanding of the unknown environment. To achieve that, we introduce curiosity-guided exploration with \textit{environment profiles}. The agent leverages these profiles to direct its curiosity-driven behaviors towards promising regions of the environment, explore diverse actions and states, and uncover the latent intentions in them.


\paragraph{Environment Profiles as Action Priors} 
Complex environments that agents are often involved in, e.g., browsers and editors, contain a wide variety of actions and states. However, users typically interact with only a subset that aligns with their specific tasks or preferences. To effectively acquire the relevant abilities to handle these problems, we capture such regularities through \textit{environment profiles}, which summarize entities, their attributes, and supported operations at the desired level of granularity.


\begin{center}
\begin{tcolorbox}[
  title=An example of \textit{Environment Profile},
  width=16cm,
  boxrule=0.45mm,
  fonttitle=\bfseries,
  enhanced,
  sharp corners,
  colback=white,
  colframe=black!70,
  arc=0mm,
  pad at break=2mm,
  outer arc=0mm
]
\texttt{%
\textbf{Entity: Map} \\
a game map with navigable roads and several points of interest, including a hospital, school, factory, home (or houses), and traffic lights.
\\
\textbf{Attributes:} \\
- road: a traversable pathway. \\
- architecture: a structure (e.g., hospital, school, home, or factory). \\
- traffic\_light: a signal requiring agents to wait before crossing. \\
\textbf{Operation:} \\
- move: move to an adjacent location.  \\
- wait\_and\_cross: wait for the signal to change and cross the intersection. \\
- enter: enter the nearest building.
}
\end{tcolorbox}
\captionof{figure}{An example of \emph{environment profile} of the sandbox in Figure \ref{fig:overview-cstb}. The elements in the map are represented by entities and attributes, while the operations enumerate possible conceptual actions for agents to build a fundamental understanding.} \label{fig:environment-profile}
\end{center}

During exploration, profiles serve as part of the initial state $s_0$ and help the LLM focus its curiosity and avoid aimless walking.

\paragraph{LLM-Driven Stochastic Exploration} We employ a high-temperature LLM as the primary source of curiosity. The elevated temperature encourages the agent to discover unconventional interaction patterns and hidden environment states, generating diverse and creative actions. At each step of exploration, the LLM observes the environment, thinks for the goals and options, and outputs an action to interact with the environment.

Formally, given the initial state $s_0$. At each step $t$, the agent samples an action from a profile-guided, high-temperature policy and then transitions according to environment dynamics:
\begin{equation}
a_t \sim \pi_\mathrm{explore}(\cdot\mid s_t,s_{t-1},s_{t-2},\dots,s_{0}),\qquad 
s_{t+1}\sim \mathcal{P}(\cdot\mid s_t,a_t),
\end{equation}
yielding a trajectory $\tau=(s_0,a_0,\ldots,s_T)$ with distribution $\rho(\tau)$.



During the exploration, we implement a two-phase exploration strategy to balance the efficiency with step allocation and myopic decision. During the initial $N_b$ steps, the agent conducts breadth-first exploration to establish a naive understanding of the environment's basic semantics, ensuring coverage of diverse regions of the action-state space. Subsequently, the agent transitions to depth-first exploration for more focused investigation of promising trajectories discovered earlier.
In this second phase, we further prevent premature convergence to a single behavior pattern by enforcing a myopic decision-making rule. Concretely, the agent only considers the most recent $N_d$ observations in $s_{t},s_{t-1},\dots,s_{t-N_d+1}$.

Through the aforementioned mechanism, we establish a transformation from the environment $\mathcal E$ to a distribution over trajectories. This process is guided by the LLM-driven stochastic policy $\pi_{\mathrm{explore}}$, which can be formalized as
\begin{equation}
    \Phi: \mathcal{E}\times\pi\times\mathcal S \to \mathcal{T}, \quad \text{where} \quad \rho = \Phi(\mathcal{E},\pi_\mathrm{explore},s_0),
\end{equation}
where $\mathcal{T}$ denotes the set of trajectories, $\rho(\tau)$ is explored with $\pi_{\mathrm{explore}}$ in $\mathcal E$, and $s_0$ concludes the environment profile.


\subsection{Adaptive Task Synthesis} \label{sec:adaptive-task-synthesis}

Having understood the environment through curiosity-guided exploration, in this section, the next step involves synthesizing training tasks that leverage the discovered environment capabilities. To achieve that, we propose adaptive task synthesis, which transforms the explored trajectories into tasks by combining the trajectory with user preferences to generate potential corresponding queries, and generating a reference solution based on the action-observation pairs.


\paragraph{Preference Guided Task Synthesis}

A key challenge in task synthesis is ensuring alignment with user preferences. To this end, we utilize \textit{user preferences} to provide structured guidance, constraining the task domain along two primary axes: task difficulty and task style.
Task difficulty is quantified by the number of entities, attributes, and operations involved in potential solutions. We adaptively control this by hierarchically scaling these factors. For instance, operations requiring coordination across multiple attributes are considered complex, while those involving few are simple. In addition to difficulty, our pipeline accepts user-defined rubrics to specify stylistic preferences, ensuring tasks reflect user demands.

This synthesis process initiates by processing raw exploration trajectories $\rho$. These trajectories are first distilled to extract salient interaction components: user inputs, LLM-generated actions, and their execution results. This distilled information is then augmented with the user preferences $u$ and provided to an LLM, which analyzes the context to synthesize a candidate query $g$ aligned with the desired constraints.

Formally, the transformation from experience to structured tasks through a preference-guided synthesis function is
\begin{equation}
    \Psi: \mathcal{T} \times \mathcal{U} \to \mathcal{G}, \quad \text{where} \quad g = \Psi\left(\rho, u\right) \sim p_\mathrm{train},
\end{equation}
where $\mathcal{U}$ denotes the space of user preferences $u$. By composing $\Phi$ and $\Psi$, we transform the environment $\mathcal E$ into its corresponding tasks $\mathcal G$ and implement $F_{\mathrm{task}}(\mathcal E) = \Psi(\Phi(\mathcal E),u)$, which yields the training task $g$ that satisfies distribution $p_{\mathrm{train}}$.

\paragraph{Reference Solution Extraction}

An essential part of our task synthesis is the extraction of reliable reference solutions. For problems posed in an unexplored environment, the ability of the agent would determine whether it could obtain a solution. In our \texttt{self-questioning} mechanism, however, problems are generated after exploration. This decoupling implies that an agent is presented with tasks whose solutions lie well beyond its current capability, and the correct answers are discoverable through the prior exploration phase. We leverage this property by designating these solutions as proxy ground-truth, i.e. reference solutions, which is shown in Figure~\ref{fig:overview-cstb}.
First, the explored trajectory is condensed into a simplified action-observation format. Then, this simplified trajectory and its corresponding task are provided to a prompted LLM to extract the complete solution process, ultimately yielding the reference solution. This solution encompasses all key steps required to solve the task.

By comparing the agent's predictions against the reference, we can robustly check correctness and provide targeted supervision.
In this way, we produce tasks of controlled difficulty and style, and ensure the availability of trustworthy solutions.


\subsection{Task Curation and Distributional Hybrid} \label{sec:task-curation}

Since $F_{\mathrm{task}}(\mathcal{E})$ produces a proxy task distribution $p_{\mathrm{train}}(g)$, the resulting samples inevitably vary in quality and difficulty. A systematic filtering mechanism is required to refine this distribution and ensure that $p_{\mathrm{train}}$ better emulates the unknown $p_{\mathrm{target}}$. To address the challenges, we propose a multi-layer filtering pipeline that evaluates the tasks across multiple dimensions during and after synthesis.

\paragraph{Real-time Filtering} During the task synthesis, the real-time filtering is applied to prevent the low-quality tasks. It operates on two primary criteria to avoid duplications. We first employ lexical overlap-based deduplication to eliminate redundant tasks. We compute token-level similarity scores between newly generated tasks and existing ones. Tasks exceeding a predefined threshold are immediately discarded in the early stage. The lexical filtering provides a computationally efficient first line of defense against redundancy. Next, semantic similarity-based check is introduced to further steer from duplicate tasks. Using embedding models, we maintain the representation of all existing tasks and guide the generation process to avoid similar trajectories.

\paragraph{Post-generation Filtering} After the initial synthesis and real-time filtering, we apply the post-filtering to ensure task quality and feasibility. This stage employs a more sophisticated evaluation. The post-filtering begins with the lexical deduplication, similar to the one in real-time filtering. Subsequently, we conduct a feasibility assessment to strictly verify task executability within the target environment. This evaluation leverages reference solutions obtained during the task synthesis phase and attempts to execute the solutions against the environment. By identifying real execution failures, potential mismatches between task requirements and environment, we effectively flag and remove synthetic tasks where the reference solution is likely hallucinatory.

\paragraph{Distributional Hybrid (Optional)}
In the default setting of interest where no external data are available, the agent relies entirely on proxy tasks, i.e.,
\begin{equation}
    p_{\mathrm{train}}(g) = F_{\mathrm{task}}(\mathcal{E})(g).
\end{equation}
Considering the superiority of mixed data in maintaining stable capabilities and expanding performance boundaries, if samples from the true target distribution $p_{\mathrm{target}}$ are available, we instead consider a hybrid distribution:
\begin{equation}
    p_{\mathrm{hybrid}}(g) = (1-\lambda) p_{\mathrm{target}}(g) + \lambda p_{\mathrm{task}}(g),
\end{equation}
with $\lambda \in (0,1]$ controlling the contribution of proxy tasks.  



\subsection{Synthetic Task Reward with LLM Judge}\label{sec:self-questioning-llm-judge}

The synthesis of tasks serves as the initial step in our \texttt{self-questioning} module. Following this, we immediately face the second challenge: reliable task reward, which is necessary for guiding agent improvement. The difficulty is further acute in complex scenarios due to the lack of ground-truth rewards.
We address this challenge by developing a reward module and providing our general reference-based LLM judge for common automated assessment of task completion quality.


\paragraph{LLM as a Judge} The LLM Judge is a broadly validated solution for self-evaluation; however, designing one for long interaction trajectories remains a challenge. To accommodate developers, we offer an LLM-based judge as a fallback proxy for $F_{\mathrm{reward}}$, providing general-purpose assessment across diverse environments without specific optimizations but with agentic abilities. Our reward module ensures both generality and practicality for users to adopt their own environment rewards.


\paragraph{Basic Principle} To better assess a long-term trajectory, our LLM judge follows two principle to carefully assign rewards:

i) \textit{Relevance and Repetition Check} serves as the primary guard, identifying trajectories that contain steps unrelated to the target task, and all hallucinated or redundant steps. Agents that exhibit those behaviors without clear justification receive zero scores immediately. 

ii) \textit{Continuous Scoring} provides nuanced assessment beyond binary success metrics. Correct solutions receive scores in high range, while incorrect attempts fall within low range, with the specific score determined by factors such as step efficiency and intermediate errors. This graduated scoring captures partial progress and provides informative feedback for agent improvement.

\paragraph{Reference-based Correctness Check} Rather than relying solely on LLM knowledge to assess agents, our judge leverages reference solutions obtained in the exploration phase. It compares agent trajectories against reference solutions to verify coverage of essential steps. The evaluation checks whether all critical steps in the reference solutions appear in the agent's execution, either directly or through functionally equivalent alternatives.

This method allows us to maintain an objective standard for rewarding. Experimental validation of the effectiveness is presented in Section \ref{sec:experiment_self-questionin}.


